\begin{pastexam}[label={px:ntu-110-dsaa-06}]{臺大資工所 110 學年度 第 VI 題}
（單選；每格答對 5 分，答錯倒扣 2.5 分）

There are 10 hexadecimal data, 53, 3B, 66, 43, 60, 5B, 7C, 14, 30, 37, inserted in the given
order into an empty hash table. The table is implemented using a circular array of 10 slots, and
one slot can only have one item. The hash function for the table is $h(k)=k \mod 9$.
\begin{enumerate}[label=(\alph*)]
\item Which data is the first one that occurs a collision? \underline{\makebox[5em]{(11)}}

(A) 43 \qquad (B) 5B \qquad (C) 14 \qquad (D) 37
\item If the hash table uses linear probing to resolve collisions and has no resizing
mechanism, which position in the array is 37 inserted? \underline{\makebox[5em]{(12)}}

(A) 3 \qquad (B) 8 \qquad (C) 4 \qquad (D) 0
\end{enumerate}
\end{pastexam}

\begin{pastsoln}
答案：(11) 選 (C) 14；(12) 選 (D) 0。

資料是十六進位，先換成十進位再算 $k \bmod 9$。 表有 $10$ 格（編號 $0$--$9$），但 $h(k)$ 只會是
$0$--$8$，第 $9$ 格只有探測時才會用到。 線性探測從 $h(k)$ 開始往後一格一格找空位，超過第 $9$ 格
繞回第 $0$ 格（資結書第~12.3~節）。

\begin{center}
\begin{tabular}{cccl}
\toprule
資料 & 十進位 & $h(k) = k \bmod 9$ & 探測的格子 $\to$ 放在 \\
\midrule
53 & $5 \cdot 16 + 3 = 83$ & $83 - 81 = 2$ & $2 \to 2$ \\
3B & $3 \cdot 16 + 11 = 59$ & $59 - 54 = 5$ & $5 \to 5$ \\
66 & $6 \cdot 16 + 6 = 102$ & $102 - 99 = 3$ & $3 \to 3$ \\
43 & $4 \cdot 16 + 3 = 67$ & $67 - 63 = 4$ & $4 \to 4$ \\
60 & $6 \cdot 16 + 0 = 96$ & $96 - 90 = 6$ & $6 \to 6$ \\
5B & $5 \cdot 16 + 11 = 91$ & $91 - 90 = 1$ & $1 \to 1$ \\
7C & $7 \cdot 16 + 12 = 124$ & $124 - 117 = 7$ & $7 \to 7$ \\
14 & $1 \cdot 16 + 4 = 20$ & $20 - 18 = 2$ & $2, 3, 4, 5, 6, 7, 8 \to 8$ \\
30 & $3 \cdot 16 + 0 = 48$ & $48 - 45 = 3$ & $3, 4, 5, 6, 7, 8, 9 \to 9$ \\
37 & $3 \cdot 16 + 7 = 55$ & $55 - 54 = 1$ & $1, 2, \ldots, 9, 0 \to 0$ \\
\bottomrule
\end{tabular}
\end{center}

\textbf{(11)} 前七筆的 $h(k)$ 是 $2, 5, 3, 4, 6, 1, 7$，互不相同，都直接放進空格。 第八筆 14 的
$h = 2$，第 $2$ 格已經有 53，是第一次碰撞，選 (C)。 (A) 43 的 $h = 4$、(B) 5B 的 $h = 1$，放進去
的時候那一格都是空的；(D) 37 也會碰撞，但在 14 之後。

\textbf{(12)} 放 37 之前，第 $1$--$9$ 格全滿（14 被推到第 $8$ 格，30 從第 $3$ 格一路推到第 $9$ 格）。
37 的 $h = 1$，從第 $1$ 格往後探測第 $1, 2, \ldots, 9$ 格都滿，繞回第 $0$ 格才是空的，放在第 $0$ 格，
選 (D)。 最後的表：
\begin{center}
\begin{tabular}{*{10}{c}}
\toprule
0 & 1 & 2 & 3 & 4 & 5 & 6 & 7 & 8 & 9 \\
\midrule
37 & 5B & 53 & 66 & 43 & 3B & 60 & 7C & 14 & 30 \\
\bottomrule
\end{tabular}
\end{center}
(A) 3、(C) 4 是 66 和 43 的格子；(B) 8 是 14 的格子。 題目說「circular array」，就是要考這一步
繞回第 $0$ 格。

\end{pastsoln}
